% !TEX root = ../main.tex

\chapter{Metodo di ricerca}
\label{cap:metodo}

\section{Obiettivi e domande di ricerca}
\label{sec:research-questions}

La rassegna condotta nel capitolo precedente restituisce un quadro
disomogeneo. Alcune categorie dell'OWASP API Security Top 10 dispongono di
tecniche di rilevazione consolidate e ampiamente validate, mentre altre
risultano trattate quasi esclusivamente in chiave dinamica o di runtime, e
altre ancora non trovano una formulazione operativa che ne consenta il
riconoscimento automatico. Questo divario non significa che alcune
categorie siano meno importanti o meno studiate, ma dipende da come si
manifesta la violazione di sicurezza in ciascuna di esse.

Da un lato, vi sono violazioni che lasciano una traccia strutturale
evidente e possono essere individuate direttamente dal codice sorgente,
dalle configurazioni o dalle specifiche del sistema. Dall'altro, vi sono
violazioni il cui riconoscimento presuppone di sapere quale utilizzo
dell'applicazione sia da considerarsi legittimo, informazione che non è
deducibile dal solo codice sorgente.

Questa frattura riflette la distinzione, ampiamente consolidata nella letteratura
sulla sicurezza del software, tra \emph{vulnerabilità tecniche} e
\emph{vulnerabilità di logica di business} (\emph{business logic vulnerabilities}).
Mentre le prime violano proprietà strutturali direttamente verificabili negli
artefatti di progetto (codice applicativo, definizioni infrastrutturali IaC o
specifiche dell'interfaccia), le seconde violano l'intento per cui il flusso
applicativo è stato concepito. Come evidenziato dalla stessa OWASP nella
motivazione di categorie quali API6:2023, le problematiche di business logic
resistono al rilevamento automatico proprio perché la correttezza è definita
rispetto all'intento del dominio applicativo.

In questo quadro, il contributo del presente lavoro non consiste nell'introdurre
una nuova categoria teorica, bensì nell'applicare sistematicamente tale
distinzione consolidata alle dieci classi dell'OWASP API Security Top 10 (2023),
verificandone le conseguenze implementative attraverso un'analisi multi-layer
degli artefatti di progetto, integrata da evidenza dinamica ove la sola
ispezione statica non sia sufficiente.

Su questa base si articola la domanda di ricerca principale del lavoro.

\rqbox{1}{Quali categorie dell'OWASP API Security Top 10 (2023)
sono riconducibili a vulnerabilità tecniche rilevabili mediante analisi
statica multi-layer su codice applicativo, definizioni infrastrutturali e
specifica dell'interfaccia, quali richiedono in aggiunta evidenza dinamica,
e quali ricadono nella logica di business richiedendo informazioni di
dominio non ricostruibili dagli artefatti analizzati?}

La domanda si scompone in due interrogativi subordinati, che ne precisano
rispettivamente la portata costruttiva e quella critica.

\rqbox{1.1}{Per le categorie riconducibili a vulnerabilità tecniche,
quale combinazione di livelli di analisi è necessaria alla loro copertura,
e quali condizioni di rischio emergono soltanto dalla congiunzione di
segnali che, presi isolatamente, non costituiscono una violazione?}

\rqbox{1.2}{Per le categorie la cui proprietà violata non è interamente
esprimibile sugli artefatti, quale componente resta comunque riconoscibile
staticamente, quale parte del residuo è accertabile mediante evidenza
dinamica e quale richiede irriducibilmente conoscenza del dominio?}

L'evidenza a sostegno di queste domande proviene da tre fonti distinte. La
mappatura per categoria presentata nel capitolo precedente fornisce la base
documentale sulla quale la distinzione viene applicata; l'implementazione dei
moduli di rilevazione e la loro validazione su fixture controllate e su un
bersaglio realistico ne verificano la traducibilità operativa; l'esclusione
motivata delle categorie inerenti alla logica di business costituisce essa stessa
un risultato, in quanto delimita con precisione il perimetro dell'approccio
anziché lasciarne i limiti impliciti. La validazione riguarda tanto i moduli
di analisi statica quanto le componenti dinamiche introdotte per le categorie
il cui residuo lo richiede.

\section{Impostazione dello studio}
\label{sec:impostazione}

Le domande formulate nella sezione precedente non ammettono una risposta
puramente documentale. Stabilire quali categorie siano riconducibili a
violazioni di natura tecnica richiede un criterio applicato in modo
sistematico, ma verificare che tale riconduzione regga sul piano operativo
richiede di costruire gli strumenti di rilevazione corrispondenti e di
osservarne il comportamento su codice reale. Lo studio assume perciò una
forma costruttiva: la risposta passa attraverso la realizzazione di un
artefatto software e la sua applicazione a bersagli definiti.

Da questa impostazione discende una precisazione necessaria sul ruolo del
sistema descritto nei capitoli seguenti. Esso non costituisce l'oggetto
della ricerca, bensì lo strumento attraverso cui la classificazione teorica
viene messa alla prova. Le scelte architetturali e implementative sono
interamente guidate dall'obiettivo metodologico: rendere operativi e
verificabili sul campo i riscontri previsti dalla classificazione. Se una
categoria è stata ricondotta a una violazione strutturale, deve esistere
una regola che la riconosca; se una categoria è stata esclusa, l'assenza del
modulo corrispondente non è una lacuna implementativa ma la conseguenza
attesa dell'analisi.

Coerentemente, si considerano evidenza a sostegno delle domande di ricerca
tre elementi soltanto: la classificazione delle dieci categorie secondo il
criterio dichiarato, l'esito verificato dell'esecuzione dei moduli di
rilevazione su artefatti di prova e su un bersaglio realistico, e la
documentazione motivata delle categorie non coperte. Non costituiscono
evidenza le dichiarazioni di copertura non accompagnate da esito di
esecuzione, né la presenza di regole nel codice non verificate su casi
concreti.

Il resto del capitolo procede in tre passi. La sezione successiva applica il
criterio alle dieci categorie e ne produce la classificazione, rispondendo
alla parte dichiarativa di RQ1. La sezione seguente deriva da tale
classificazione i requisiti cui uno strumento di rilevazione deve
soddisfare, senza ancora impegnarsi su una realizzazione particolare. Il
capitolo si chiude con il disegno della valutazione, ossia i bersagli, la
costruzione del riferimento di verità e le grandezze osservate.

\section{Applicazione del criterio alle dieci categorie}
\label{sec:applicazione-criterio}

La classificazione muove dalla distinzione introdotta nella
Sezione~\ref{sec:research-questions} e la applica a ciascuna delle dieci
categorie dell'OWASP API Security Top 10 (2023) secondo una procedura
uniforme. Per ogni categoria si individua anzitutto la proprietà di
sicurezza la cui violazione la definisce, attenendosi alla formulazione
ufficiale e alla caratterizzazione emersa dalla rassegna del capitolo
precedente. Si valuta quindi se tale proprietà sia esprimibile come
predicato sugli artefatti di progetto, ossia se esista un tratto del codice
applicativo, delle definizioni infrastrutturali o della specifica
dell'interfaccia la cui presenza o assenza sia sufficiente a stabilire la
violazione.

L'applicazione del criterio produce tre esiti distinti. Una categoria è
detta \emph{tecnica} quando la proprietà violata è interamente esprimibile
in tali termini: la valutazione del predicato richiede il solo esame degli
artefatti. È detta \emph{di logica di business} quando la proprietà violata
è definita per riferimento all'intento applicativo, e nessun esame degli
artefatti può stabilirne la violazione in assenza di informazioni sul
comportamento atteso. È detta infine \emph{ibrida} quando la proprietà si
decompone in una componente strutturale e in un residuo di dominio: in
questi casi l'analisi statica riconosce una condizione necessaria alla
violazione, ma non sufficiente a stabilirla.

Il residuo che ne risulta non è omogeneo. In alcuni casi esso è accertabile
osservando il comportamento del sistema in esecuzione, poiché la proprietà
mancante si manifesta nella risposta a una sollecitazione mirata; in altri
richiede una conoscenza dell'intento applicativo che nessuna osservazione
del sistema, statica o dinamica, può restituire. La distinzione fra i due
tipi di residuo separa le categorie ibride da quelle di logica di business.

La classe intermedia non rappresenta un'incertezza della classificazione, ma
una proprietà del materiale classificato. Diverse categorie descrivono
violazioni di controlli di autorizzazione la cui assenza è strutturalmente
riconoscibile, mentre la determinazione di quali soggetti debbano accedere a
quali risorse appartiene al dominio applicativo. Trattare tali categorie
come interamente tecniche sovrastimerebbe la portata dell'analisi statica;
trattarle come interamente di logica di business ne trascurerebbe la
componente effettivamente rilevabile. La loro individuazione costituisce
perciò l'oggetto proprio di RQ1.2, che interroga la frazione di proprietà
approssimabile staticamente e il residuo che non lo è.

La Tabella~\ref{tab:classificazione} riporta l'esito dell'applicazione del
criterio; i paragrafi che seguono si limitano a motivare i casi che
richiedono una precisazione.

\begin{table}[htbp]
  \centering
  \small
  \linespread{1.1}\selectfont
  \begin{tabularx}{\textwidth}{@{}l l >{\raggedright\arraybackslash}X >{\raggedright\arraybackslash}X@{}}
    \toprule
    \textbf{Categoria} & \textbf{Classe} & \textbf{Componente strutturale} & \textbf{Residuo} \\
    \midrule
    API7 & tecnica & flusso da ingresso non fidato a chiamata in uscita senza validazione & --- \\
    API8 & tecnica & difformità da una baseline nota & quale baseline sia corretta per il contesto d'uso \\
    \addlinespace
    API1 & ibrida & assenza di controllo di autorizzazione sul percorso verso l'oggetto & chi possa accedere a quale oggetto (dinamico) \\
    API3 & ibrida & flusso di dati fra richiesta e risposta & quali proprietà siano sensibili (dinamico) \\
    API5 & ibrida & funzione riservata raggiungibile senza controllo & ripartizione dei privilegi prevista (dinamico) \\
    API4 & ibrida & assenza di un limite a un dato livello & se il limite sia imposto altrove nel deployment \\
    API2 & ibrida & gestione dei token e scelte crittografiche & abuso dei flussi di autenticazione (dinamico) \\
    API10 & ibrida & canale non cifrato, assenza di limiti sulle risorse acquisite & contratto del fornitore \\
    \addlinespace
    API6 & logica di business & --- & quali flussi siano sensibili \\
    \addlinespace
    API9 & fuori griglia & discrepanza fra contratto e codice & sistemi ignoti all'organizzazione \\
    \bottomrule
  \end{tabularx}
  \vspace{0.6em}
  \caption{Esito dell'applicazione del criterio alle dieci categorie
    dell'OWASP API Security Top 10 (2023). Il residuo è marcato come
    \emph{dinamico} quando è accertabile sollecitando il sistema in
    esecuzione.}
  \label{tab:classificazione}
\end{table}

Fra le categorie tecniche, API7 soddisfa il criterio senza riserve: il
difetto è nella mancata validazione dell'indirizzo, non nell'uso che
l'applicazione fa della risorsa recuperata. API8 richiede invece una
precisazione: il confronto con la baseline è meccanico, ma la baseline
stessa può dipendere dall'uso previsto del sistema. Una politica di
condivisione fra origini che ammetta qualunque chiamante è un difetto per
un'interfaccia a client autenticati e una scelta legittima per
un'interfaccia pubblica di sola lettura. La componente di dominio non
riguarda quindi la rilevazione della configurazione, ma la determinazione
di quale configurazione sia corretta.

Le tre categorie di autorizzazione (API1, API3, API5) condividono la stessa
struttura a granularità diversa --- oggetti, proprietà, funzioni --- ed è
questa condivisione a rendere la classe ibrida un tratto del materiale e
non un artificio: l'assenza del controllo è riconoscibile sugli artefatti,
mentre la politica che il controllo dovrebbe applicare risiede nel modello
di autorizzazione e si accerta sollecitando l'interfaccia con identità
distinte. Le altre tre ibride hanno residui di natura diversa. Per API4 il
residuo non riguarda l'intento applicativo ma il contesto di deployment,
ed è colmabile estendendo l'analisi alle definizioni infrastrutturali; per
API2 la linea di separazione è quella già individuata dalla rassegna, fra
proprietà deterministiche dei token e abusi dei flussi osservabili solo in
esecuzione; per API10 le cinque condizioni che la compongono si
distribuiscono fra le due componenti.

API6 è l'unica categoria priva di componente strutturale: due applicazioni
sintatticamente identiche possono differire quanto alla presenza della
vulnerabilità a seconda del modello di business che le sostiene, e il
residuo non è accertabile nemmeno per via dinamica, poiché l'osservazione
restituisce ciò che il sistema fa, non ciò che dovrebbe consentire.

API9, infine, eccede la griglia perché il criterio non le si applica: le
altre categorie descrivono la violazione di una proprietà di un sistema
noto, API9 l'esistenza di sistemi che l'organizzazione ignora di
possedere. Il confronto fra ciò che il contratto dichiara e ciò che il
codice realizza ne individua una parte, ma quanto sfugge a entrambe le
rappresentazioni resta fuori portata per costruzione.

L'esito è fortemente asimmetrico: due categorie tecniche, sei ibride con
residui di entità variabile, una di logica di business e una fuori
classificazione. Questo costituisce di per sé una risposta alla parte
dichiarativa di RQ1, e ridimensiona l'aspettativa che l'analisi degli
artefatti di progetto possa coprire in modo uniforme il perimetro
descritto dall'OWASP API Security Top 10.

\section{Requisiti derivati per lo strumento di rilevazione}
\label{sec:requisiti}

La classificazione appena prodotta non descrive uno strumento, ma vincola
ciò che uno strumento adeguato deve saper fare. I requisiti che seguono
sono formulati in termini di capacità richieste, indipendentemente dalle
tecnologie che potranno realizzarle; la loro traduzione in un'architettura
concreta è oggetto del capitolo successivo.

Il primo requisito riguarda l'eterogeneità degli artefatti. Le proprietà
individuate dalla classificazione non risiedono tutte nello stesso luogo:
alcune si manifestano nel codice applicativo, altre nelle definizioni
infrastrutturali, altre ancora nella specifica dell'interfaccia. Uno
strumento che esamini un solo tipo di artefatto è quindi strutturalmente
incapace di coprire il perimetro, non per limiti di implementazione ma per
la natura delle proprietà in gioco. Ne discende la necessità di un'analisi
articolata su livelli distinti, ciascuno dotato di una rappresentazione
adeguata al tipo di artefatto che esamina.

Il secondo requisito riguarda la relazione fra i livelli. Alcune delle
condizioni emerse dalla classificazione non sono osservabili in alcun
singolo artefatto, ma nella discrepanza fra due di essi. Un percorso
presente nel codice e assente dal contratto non costituisce una violazione
in nessuno dei due luoghi presi separatamente, eppure la loro
divergenza segnala una superficie esposta e non governata. Analogamente,
l'assenza di un limite nel codice applicativo acquista significato diverso
a seconda che le definizioni infrastrutturali ne prevedano uno altrove.
Lo strumento deve dunque disporre di un meccanismo che metta in relazione
segnali provenienti da livelli diversi e ne derivi giudizi che nessun
livello esprime autonomamente. Questo requisito risponde direttamente a
RQ1.1.

Il terzo requisito riguarda il residuo delle categorie ibride. Per quelle
in cui esso è accertabile osservando il comportamento del sistema in
esecuzione, lo strumento deve poter sollecitare l'interfaccia e valutarne
le risposte, poiché la componente strutturale da sola non consente di
stabilire la violazione ma soltanto la sua possibilità. L'integrazione di
evidenza dinamica non costituisce quindi un'estensione opzionale, bensì
la conseguenza diretta della classificazione: senza di essa le categorie
ibride resterebbero indecise.

Il quarto requisito riguarda la tracciabilità. Poiché lo strumento serve a
verificare una classificazione, ogni segnalazione da esso prodotta deve
essere riconducibile alla categoria di cui costituisce evidenza e al
livello di analisi che l'ha generata. Una segnalazione priva di tale
riconduzione può indicare un difetto reale, ma non contribuisce a
rispondere alle domande di ricerca, poiché non è imputabile ad alcuna
delle classi individuate.

\section{Disegno della valutazione}
\label{sec:disegno-valutazione}

La valutazione ha due scopi: controllare che ogni modulo di rilevazione
riconosca davvero la condizione per cui è stato scritto, e osservare come
si comporta lo strumento nel suo insieme su codice che non è stato scritto
per essere analizzato. Al primo scopo servono coppie di artefatti di prova,
ciascuna formata da una versione che contiene la condizione da rilevare e
da una che non la contiene: il modulo deve segnalare la prima e tacere
sulla seconda. Poiché artefatti e regole hanno lo stesso autore, questa
prova dice solo che l'implementazione fa ciò che si intendeva farle fare,
non che funzioni su codice reale. Al secondo scopo serve OWASP crAPI,
un'applicazione volutamente vulnerabile usata come riferimento in
letteratura e qui fissata a una revisione precisa per rendere i risultati
riproducibili. Non essendo scritta dall'autore, riduce la circolarità della
prima prova; resta però un bersaglio didattico, costruito perché le sue
vulnerabilità siano riconoscibili, e quindi i risultati non si possono
estendere senza cautela ad applicazioni reali. Per gli artefatti di prova
il riferimento di verità è noto per costruzione; per crAPI è stato
ricostruito a mano a partire dalle vulnerabilità documentate dal progetto,
ed è riportato per esteso così che i risultati possano essere giudicati
rispetto a ciò che si è effettivamente assunto come vero. Due segnalazioni
si considerano equivalenti quando riguardano la stessa categoria OWASP e lo
stesso punto di esposizione dell'interfaccia, identificato da metodo e
percorso normalizzato, a prescindere dal livello di analisi che le ha
prodotte; la convenzione è stata fissata prima di osservare i risultati,
per non scegliere a posteriori la granularità più conveniente. Per ogni
categoria si registra se i moduli segnalano gli artefatti che contengono la
condizione e non quelli che ne sono privi; le categorie senza alcuna
segnalazione sono distinte fra quelle in cui ciò discende dalla
classificazione e quelle in cui indica una copertura mancante, e la sola
presenza di una regola nel codice, senza l'esito della sua esecuzione, non
conta come copertura. Va infine detto che la classificazione di questo
capitolo è stata formalizzata quando parte dei moduli era già stata
realizzata: il criterio deriva dalla rassegna della letteratura e non dalle
capacità dello strumento, ma l'ordine reale del lavoro non coincide con
quello dell'esposizione, e i risultati vanno letti tenendone conto.
